\subsection{Методика экспериментов}
\label{subsec:methodology}

Проведены две серии экспериментов, соответствующие двум основным сценариям
использования Apache Paimon: пакетная аналитика над таблицами только для добавления и
near-real-time-витрина на таблице с первичным ключом при потоковой записи.

\subsubsection{Сценарий 1: пакетная аналитика на наборе TPC-DS}

TPC-DS~\cite{tpcds-spec} --- стандартизированный набор для оценки систем поддержки принятия решений:
схема розничной торговли (таблицы фактов \texttt{store\_sales}, \texttt{web\_sales},
\texttt{catalog\_sales} и связанные измерения), генератор данных с управляемым
масштабом и \num{99} шаблонов запросов (с учётом вариантов --- \num{103} запроса),
покрывающих сканирование, фильтрацию, соединения, группировку и оконные функции.
Использован масштаб \SI{100}{\giga\byte}; данные сгенерированы, загружены в Apache
Paimon как таблицы только для добавления, разбиты на разделы (типичный размер файла
раздела --- порядка \SI{256}{\mega\byte}), параллелизм запросов --- 300. Сравнивались
форматы Parquet и Vortex по времени выполнения каждого запроса; для каждого запроса
выполнено несколько повторов, бралось характерное (устойчивое) значение. Запись в
этом сценарии не измерялась: загрузка крупных наборов упиралась в пропускную
способность объектного хранилища, что не позволяло корректно сопоставить форматы по
записи (этот сценарий покрывается второй серией).

\subsubsection{Сценарий 2: near-real-time-витрина при потоковой записи}

Разработано задание Flink, моделирующее near-real-time-витрину. Используется таблица
с первичным ключом, полем последовательности и полем вида изменения (схема захвата
изменений). Схема таблицы --- журнал событий из 25 столбцов, подобранный так, чтобы
дать форматам возможность проявить разнообразие кодеков: категориальные строки
низкой кардинальности (страна, тип устройства, операционная система, метод запроса
--- словарное и RLE-кодирование), строки с регулярной структурой и общими
префиксами (идентификаторы сессии и запроса, путь URL, версия ОС, строка
user-agent --- FSST), монотонные и узкодиапазонные целые (код ответа, задержка,
ширина экрана --- FOR и бит-упаковка), числа с плавающей точкой (координаты, выручка
--- ALP), столбцы с большой долей null (идентификатор родительского запроса ---
\SI{30}{\percent}, идентификатор объявления --- \SI{50}{\percent}, выручка ---
\SI{90}{\percent}). Ширина в 25 столбцов сопоставима с таблицами фактов TPC-DS ---
реалистичная форма аналитической таблицы. Состав схемы приведён в
таблице~\ref{tab:schema}.

\begin{table}
    \caption{Схема таблицы near-real-time-витрины (25 столбцов)}
    \label{tab:schema}
    {\small
    \begin{tabularx}{\textwidth}{|>{\raggedright\arraybackslash}p{0.22\textwidth}|>{\raggedright\arraybackslash}p{0.40\textwidth}|X|}\hline
    \textbf{Группа} & \textbf{Столбцы} & \textbf{Ожидаемые кодеки} \\ \hline
    Ключ и служебные & \texttt{id} (первичный ключ), \texttt{seq} (последовательность), \texttt{op} (вид изменения) & дельта / FOR; константа в фазе обновления \\ \hline
    Категории & \texttt{country\_code}, \texttt{device\_type}, \texttt{os\_name}, \texttt{event\_type}, \texttt{request\_method}, \texttt{referrer\_host} & словарь, RLE \\ \hline
    Префиксные строки & \texttt{session\_id}, \texttt{request\_id}, \texttt{parent\_req\_id}, \texttt{url\_path}, \texttt{user\_agent}, \texttt{os\_version} & FSST \\ \hline
    Целые & \texttt{user\_id}, \texttt{status\_code}, \texttt{latency\_ms}, \texttt{bytes\_sent}, \texttt{screen\_width}, \texttt{ad\_id} & FOR, бит-упаковка, дельта \\ \hline
    С плавающей точкой & \texttt{geo\_lat}, \texttt{geo\_lon}, \texttt{revenue} & ALP \\ \hline
    Временная метка & \texttt{ts} & дельта \\ \hline
    \end{tabularx}
    }
\end{table}

Эксперимент состоит из четырёх фаз. INIT --- начальная загрузка \num{50000000} строк
с последовательными идентификаторами (поле последовательности равно нулю); измеряется
пропускная способность записи. UPDATE --- поток обновлений по случайным
идентификаторам из загруженного диапазона (поле последовательности возрастает) с
ограничением скорости (для получения воспроизводимой нагрузки скорость зафиксирована
на уровне порядка 100\,000~строк/с; объём подобран так, чтобы фаза длилась
сопоставимое время в разных запусках). Интервал чекпойнтов Flink установлен в
240~с --- это режим near-real-time: фиксации происходят раз в несколько минут,
файлы уровня~0 крупные. DURING --- во время фазы обновления параллельно, с
интервалом в несколько минут, запускаются пакетные аналитические запросы по текущему
снимку (несколько выборок); это моделирует потребителя, читающего витрину во время её
наполнения. FINAL --- после завершения записи те же запросы выполняются по
итоговому снимку. Различие фаз DURING и FINAL принципиально: в фазе DURING читатель
конкурирует с активной записью за ввод-вывод, а LSM-дерево содержит много ещё не
слитых файлов уровня~0; в фазе FINAL запись остановлена, а дерево относительно
упорядочено (хотя полное принудительное слияние перед итоговыми запросами не
выполняется --- читается состояние «как есть после остановки записи», что и
соответствует реальной витрине).

Набор аналитических запросов охватывает характерные типы нагрузки; он приведён в
таблице~\ref{tab:queries}. Каждый запрос материализует результат в таблицу-приёмник
(чтобы исключить влияние клиента на измерение); метрикой служит время выполнения в
миллисекундах. Весь цикл (создание таблицы --- INIT --- UPDATE с DURING-выборками ---
FINAL) повторяется несколько раз (число повторов --- не менее трёх); состояние между
повторами сбрасывается. По повторам и DURING-выборкам вычисляются как среднее, так
и медиана; основной метрикой принята медиана, так как DURING-сэмплы могут
попадать в момент массового сброса данных по чекпойнту Flink (в эти моменты writer
одновременно пишет в объектное хранилище и данные таблицы, и состояние чекпойнта,
размещаемое в той же директории хранилища), что приводит к единичным «хвостовым»
сэмплам длительностью в десятки секунд; такие сэмплы инфлируют среднее, но не
влияют на медиану. Среднее приводится отдельно как индикатор подверженности
формата выбросам.

\begin{xltabular}{\textwidth}{|>{\raggedright\arraybackslash}p{0.40\textwidth}|X|}
    \caption{Аналитические запросы near-real-time-сценария}\label{tab:queries} \\ \hline
    \textbf{Запрос} & \textbf{Содержание и проверяемый механизм} \\ \hline
    \endfirsthead
    \caption*{Продолжение таблицы~\ref{tab:queries}} \\ \hline
    \textbf{Запрос} & \textbf{Содержание и проверяемый механизм} \\ \hline
    \endhead
    \hline
    \endfoot
    \hline
    \endlastfoot
    \texttt{POINT} & выборка одной строки по первичному ключу; точечный доступ, фильтры Блума и статистики \\ \hline
    \texttt{FULLSCAN\_PROJ} & полное сканирование с проекцией; нет предиката --- проверяет «чистую» скорость декодирования и накладные расходы границы \\ \hline
    \texttt{RANGE\_SESSION} & диапазонный фильтр по \texttt{session\_id} (строка-префикс); FSST + пропуск файлов по статистикам строкового столбца \\ \hline
    \texttt{RANGE\_URL} & диапазонный фильтр по \texttt{url\_path}; то же на пути URL \\ \hline
    \texttt{AGG\_SUM\_BYTES} & \texttt{SUM} по числовому столбцу с полным сканированием; агрегация без селективного предиката \\ \hline
    \texttt{AGG\_COUNT\_DISTINCT} & \texttt{COUNT(DISTINCT)} по строковому столбцу; тяжёлая агрегация с полным сканированием \\ \hline
    \texttt{AGG\_FILTERED} & \texttt{COUNT} с диапазонным предикатом; предикат отсекает данные, агрегация амортизируется \\ \hline
    \texttt{AGG\_GROUP\_COUNTRY} & \texttt{GROUP BY} по категориальному столбцу; группировка по столбцу низкой кардинальности \\ \hline
    \texttt{AGG\_MULTIDIM} & многомерная группировка с несколькими агрегатами \\ \hline
    \texttt{FUNNEL} & многошаговый запрос (последовательность событий пользователя); несколько проходов и предикатов \\ \hline
    \texttt{JOIN\_USER\_ACTIVITY} & соединение по \texttt{user\_id} с агрегацией по типу события; column pruning при соединении \\ \hline
    \texttt{JOIN\_SESSION\_LATENCY} & соединение по \texttt{session\_id} с диапазонным фильтром по задержке \\ \hline
    \texttt{JOIN\_REVENUE\_PROFILE} & подзапрос-сборка активных пользователей и соединение с их строками покупок; 5 из 25 столбцов \\ \hline
\end{xltabular}

По итогам предварительных запусков зафиксирован
рабочий набор параметров: распределение слотов между записью и чтением (общая сумма
--- весь кластер); число бакетов; размер буфера записи под выбранный интервал
чекпойнтов; увеличенный размер пакета записи Vortex (снижает накладные расходы JNI
при чтении); полные статистики по столбцам (\texttt{metadata.stats-mode = full});
расширенные политики хранения снимков и манифестов (необходимы для устойчивости при
частых фиксациях и параллельном чтении); файловые форматы --- без специфичных тонких
настроек размера файлов и страниц. От ряда гипотез пришлось отказаться: использование
строчного формата Avro для свежих файлов уровня~0 ухудшает аналитику Vortex на
итоговом снимке; агрессивные настройки размера файлов и страниц благоприятствуют
Parquet и раздувают хранилище за счёт накопления удаляемых файлов в окне хранения;
эти результаты обсуждаются в подразделе~\ref{subsec:issues}.

Измерение пропускной способности записи требует оговорки. В фазе обновления все
строки имеют одинаковое значение поля последовательности, поэтому Paimon
дедуплицирует записи с совпадающим ключом и последовательностью в буфере записи до
сброса на диск. При большом числе обновлений на один ключ значительная доля строк не
доходит до хранилища, и измеренная «пропускная способность» в фазе обновления
отражает скорость генерации и дедупликации в памяти, а не скорость записи на диск.
Поэтому в выводах по записи используется в первую очередь пропускная способность фазы
INIT (последовательные ключи, реальная запись), а показатели фазы UPDATE
интерпретируются с этой оговоркой.

Однократный запуск даёт разброс результатов до 25--40\,\% между
идентичными конфигурациями: на результат влияют разогрев JIT-компилятора первого из
прогоняемых форматов, стоимость первой загрузки нативной библиотеки, размещение
процессов по узлам, состояние пулов соединений с хранилищем, накопленное состояние
кластера. Для статистически значимого сравнения требуется несколько повторов с
чередованием порядка форматов; единичные выбросы
интерпретировались как артефакты совпадения момента запроса с пиком ввода-вывода от
записи или паузой сборки мусора.
